\documentclass[11pt,a4paper]{article}
\usepackage[utf8]{inputenc}
\usepackage[T1]{fontenc}
\usepackage[spanish,es-minimal,es-noshorthands]{babel}
\AtBeginDocument{\renewcommand{\sin}{\operatorname{sen}}\renewcommand{\arcsin}{\operatorname{arcsen}}}
\usepackage{lmodern}
\usepackage[sfdefault,lining]{FiraSans}
\usepackage[a4paper,top=2.3cm,bottom=2.2cm,left=1.8cm,right=1.8cm,headheight=1.3cm,headsep=0.4cm,footskip=1cm]{geometry}
\usepackage{amsmath,amssymb}
\usepackage{icomma}
\usepackage{graphicx}
\usepackage[table]{xcolor}
\usepackage{tikz}
\usepackage{fontawesome5}
\usepackage{ccicons}
\usepackage{fancyhdr}
\usepackage{lastpage}
\usepackage{tasks}
\usepackage{needspace}
\usepackage[most]{tcolorbox}
\usepackage{array}
\usepackage[hidelinks]{hyperref}
\AtBeginDocument{\hypersetup{pdftitle={Soluciones de la ficha de la unidad 1: Números naturales y enteros. 2º ESO},pdfsubject={Matemáticas 2º ESO},pdfkeywords={números naturales, división entera, múltiplos, divisores, números primos, mcm, MCD, números enteros, operaciones combinadas, potencias, raíces}}}

\definecolor{navy}{HTML}{1B2748}
\definecolor{accent}{HTML}{E8603A}
\definecolor{teal}{HTML}{1F9E90}
\definecolor{soft}{HTML}{EEF1F8}
\definecolor{muted}{HTML}{5D6685}

\setlength{\parindent}{0pt}
\setlength{\parskip}{3pt}
\settasks{label=\alph*), label-format=\bfseries\color{navy}, label-width=1.3em, item-indent=1.8em, column-sep=0.6em, after-item-skip=2pt, before-skip=3pt, after-skip=4pt}

% ---------- capçalera i peu ----------
\newcommand{\llicencia}{\href{https://creativecommons.org/licenses/by-nc-sa/4.0/deed.es}{\ccbyncsa}\ \ Licencia Creative Commons Atribución-NoComercial-CompartirIgual 4.0 Internacional (CC BY-NC-SA 4.0)}
\pagestyle{fancy}
\fancyhf{}
\fancyhead[L]{\raisebox{-0.35\height}{\includegraphics[height=0.95cm]{logo-gva.png}}}
\fancyhead[C]{\small\color{muted}Soluciones de la ficha · Unidad 1: Números naturales y enteros · 2\textsuperscript{o} ESO}
\fancyhead[R]{\raisebox{-0.35\height}{\includegraphics[height=0.95cm]{logo-vermellar.png}}}
\fancyfoot[L]{\scriptsize\color{muted}\llicencia}
\fancyfoot[R]{\small\color{muted}\thepage}
\renewcommand{\headrulewidth}{0.4pt}
\renewcommand{\headrule}{\color{muted!40}\hrule height \headrulewidth}
\fancypagestyle{primera}{\fancyhf{}\renewcommand{\headrulewidth}{0pt}
  \fancyfoot[L]{\scriptsize\color{muted}\llicencia}
  \fancyfoot[R]{\small\color{muted}\thepage}}

% ---------- nivells i exercicis ----------
\newcommand{\estrelles}[1]{{\color{accent}\ifcase#1\or\faStar\faStar[regular]\faStar[regular]\or\faStar\faStar\faStar[regular]\or\faStar\faStar\faStar\fi}}
\newcounter{exer}
\newcommand{\ex}[1]{\par\needspace{4\baselineskip}\medskip\stepcounter{exer}\noindent{\color{navy}\bfseries\theexer.}\ \ {\small\estrelles{#1}}\ \ }

\newcommand{\seccio}[2]{\par\needspace{10\baselineskip}\bigskip
  \begin{tcolorbox}[colback=navy,colframe=navy,coltext=white,arc=2mm,boxrule=0pt,left=3mm,right=3mm,top=1.2mm,bottom=1.2mm]
  {\large\bfseries #1\quad #2}
  \end{tcolorbox}\vspace{-2pt}}

\newcolumntype{C}[1]{>{\centering\arraybackslash}p{#1}}

% quadre de recordatori (només a la fitxa d'exercicis)
\newcommand{\recorda}[1]{\par\smallskip\begin{tcolorbox}[colback=accent!7,colframe=accent!55,arc=2mm,boxrule=0.5pt,left=2.5mm,right=2.5mm,top=1mm,bottom=1mm]\small{\color{accent}\faLightbulb[regular]}\ \textbf{Recuerda:}\ #1\end{tcolorbox}}
% recta numèrica entera: \rectabuida{min}{max} i \recta{min}{max}{punts}
\newcommand{\rectabase}[2]{\draw[latex-latex,muted] (#1-0.8,0)--(#2+0.8,0);
  \foreach \i in {#1,...,#2}{\draw[muted] (\i,0.12)--(\i,-0.12) node[below=-1pt,font=\scriptsize,text=muted]{$\i$};}}
\newcommand{\rectabuida}[2]{\begin{tikzpicture}[x=0.78cm,baseline=0]\rectabase{#1}{#2}\end{tikzpicture}}
\newcommand{\recta}[3]{\begin{tikzpicture}[x=0.78cm,baseline=0]\rectabase{#1}{#2}
  \foreach \p in {#3}{\fill[accent] (\p,0) circle (0.11) node[above=1pt,font=\small\bfseries,text=navy]{$\p$};}\end{tikzpicture}}
\newcommand{\res}[1]{\ifmmode\boxed{#1}\else\fbox{#1}\fi}

% =====================================================================
% Capçalera FentMates de les fitxes: un color per a cada curs.
% Els colors són els de la pàgina inicial (LEVELS a assets/js/materials.js);
% el logo ve de _eines/logo/pdf/ (python3 _eines/logo/genera.py).
%
% Ús, al final del preàmbul de la unitat:
%   % =====================================================================
% Capçalera FentMates de les fitxes: un color per a cada curs.
% Els colors són els de la pàgina inicial (LEVELS a assets/js/materials.js);
% el logo ve de _eines/logo/pdf/ (python3 _eines/logo/genera.py).
%
% Ús, al final del preàmbul de la unitat:
%   % =====================================================================
% Capçalera FentMates de les fitxes: un color per a cada curs.
% Els colors són els de la pàgina inicial (LEVELS a assets/js/materials.js);
% el logo ve de _eines/logo/pdf/ (python3 _eines/logo/genera.py).
%
% Ús, al final del preàmbul de la unitat:
%   \input{../comu/fentmates.tex}
%   \fmcurs{1eso}
% i, al començament del document:
%   \capcalera{Matemàtiques}{1r ESO}{Fitxa d'exercicis · Unitat 2}{Divisibilitat}
% Opcional abans de \input: \def\fmneed{14} (espai mínim abans de cada secció).
% =====================================================================
\makeatletter
\def\fm@logodir{../../logo/pdf}
\providecommand{\fmneed}{10}
\def\fm@cursnom{}

% clau del curs, color i color del text damunt del color
\newcommand{\fm@curs}[3]{\@namedef{fm@color@#1}{#2}\@namedef{fm@text@#1}{#3}}
\fm@curs{1eso}{E6B84A}{navy}
\fm@curs{2eso}{3B82F6}{white}
\fm@curs{3eso}{FF7A50}{white}
\fm@curs{4eso-a}{A78BFA}{white}
\fm@curs{4eso-b}{6D28D9}{white}
\fm@curs{1batx-mi}{9B1C1C}{white}
\fm@curs{1batx-ccss}{EF6B6B}{white}
\fm@curs{2batx-mii}{1E6B47}{white}
\fm@curs{2batx-ccss}{2BB3A3}{white}

\newcommand{\fmcurs}[1]{%
  \@ifundefined{fm@color@#1}{\PackageError{fentmates}{Curs desconegut: #1}{}}{}%
  \edef\fm@tmp{\noexpand\definecolor{curs}{HTML}{\@nameuse{fm@color@#1}}%
               \noexpand\colorlet{curstext}{\@nameuse{fm@text@#1}}}\fm@tmp
  \colorlet{cursfosc}{curs!72!black}%
  \def\fm@logo{\fm@logodir/vertical-#1.pdf}%
  \def\fm@marca{\fm@logodir/marca-#1.pdf}}

% ---------- capçalera de la primera pàgina ----------
% #1 matèria (petit), #2 curs, #3 tipus de document i unitat, #4 títol
\newcommand{\capcalera}[4]{%
  \gdef\fm@cursnom{#2}%
  \par\noindent
  \begin{tcolorbox}[colback=curs,colframe=curs,boxrule=0pt,arc=2.2mm,width=3.5cm,height=2.15cm,
      valign=center,left=3mm,right=2mm,top=1mm,bottom=1mm,nobeforeafter,box align=center]
    \color{curstext}{\footnotesize #1}\par\vspace{1pt}{\fontsize{19}{21}\selectfont\bfseries #2\par}%
  \end{tcolorbox}\hspace{5mm}%
  \begin{minipage}[c]{\dimexpr\linewidth-3.5cm-5mm-2.9cm\relax}
    {\color{cursfosc}\small\bfseries\MakeUppercase{#3}}\par\vspace{2pt}
    {\color{navy}\fontsize{20}{23}\selectfont\bfseries #4\par}
  \end{minipage}\hfill
  \begin{minipage}[c]{2.6cm}\raggedleft\includegraphics[width=2.6cm]{\fm@logo}\end{minipage}\par
  \vspace{2pt}}

% ---------- barres de secció del color del curs ----------
\renewcommand{\seccio}[2]{\par\needspace{\fmneed\baselineskip}\bigskip
  \begin{tcolorbox}[colback=curs,colframe=curs,coltext=curstext,arc=2mm,boxrule=0pt,left=3mm,right=3mm,top=1.2mm,bottom=1.2mm]
  {\large\bfseries #1\quad #2}
  \end{tcolorbox}\vspace{-2pt}}

% ---------- estrelles de dificultat del color del curs (a Batxillerat no se'n posen) ----------
\@ifundefined{estrelles}{}{%
  \renewcommand{\estrelles}[1]{{\color{curs}\ifcase#1\or\faStar\faStar[regular]\faStar[regular]\or\faStar\faStar\faStar[regular]\or\faStar\faStar\faStar\fi}}}

% ---------- capçalera i peu de les pàgines ----------
% primera pàgina: logos institucionals al peu
\fancypagestyle{primera}{\fancyhf{}\renewcommand{\headrulewidth}{0pt}
  \fancyfoot[L]{\raisebox{-0.32\height}{\includegraphics[height=0.78cm]{logo-gva.png}}\hspace{5mm}%
    \raisebox{-0.32\height}{\includegraphics[height=0.78cm]{logo-vermellar.png}}}
  \fancyfoot[R]{\scriptsize\color{muted}\href{https://creativecommons.org/licenses/by-nc-sa/4.0/deed.ca}{\ccbyncsa}\ \ CC BY-NC-SA 4.0 · fentmates.org\qquad\small\thepage}}
% la resta: el símbol a l'esquerra i el curs a la dreta, amb el títol al centre com abans
\fancyhead[L]{\raisebox{-0.3\height}{\includegraphics[height=0.62cm]{\fm@marca}}}
\fancyhead[R]{\small\bfseries\color{cursfosc}\fm@cursnom}
\renewcommand{\headrule}{\color{curs!45}\hrule height \headrulewidth}
\makeatother

%   \fmcurs{1eso}
% i, al començament del document:
%   \capcalera{Matemàtiques}{1r ESO}{Fitxa d'exercicis · Unitat 2}{Divisibilitat}
% Opcional abans de \input: \def\fmneed{14} (espai mínim abans de cada secció).
% =====================================================================
\makeatletter
\def\fm@logodir{../../logo/pdf}
\providecommand{\fmneed}{10}
\def\fm@cursnom{}

% clau del curs, color i color del text damunt del color
\newcommand{\fm@curs}[3]{\@namedef{fm@color@#1}{#2}\@namedef{fm@text@#1}{#3}}
\fm@curs{1eso}{E6B84A}{navy}
\fm@curs{2eso}{3B82F6}{white}
\fm@curs{3eso}{FF7A50}{white}
\fm@curs{4eso-a}{A78BFA}{white}
\fm@curs{4eso-b}{6D28D9}{white}
\fm@curs{1batx-mi}{9B1C1C}{white}
\fm@curs{1batx-ccss}{EF6B6B}{white}
\fm@curs{2batx-mii}{1E6B47}{white}
\fm@curs{2batx-ccss}{2BB3A3}{white}

\newcommand{\fmcurs}[1]{%
  \@ifundefined{fm@color@#1}{\PackageError{fentmates}{Curs desconegut: #1}{}}{}%
  \edef\fm@tmp{\noexpand\definecolor{curs}{HTML}{\@nameuse{fm@color@#1}}%
               \noexpand\colorlet{curstext}{\@nameuse{fm@text@#1}}}\fm@tmp
  \colorlet{cursfosc}{curs!72!black}%
  \def\fm@logo{\fm@logodir/vertical-#1.pdf}%
  \def\fm@marca{\fm@logodir/marca-#1.pdf}}

% ---------- capçalera de la primera pàgina ----------
% #1 matèria (petit), #2 curs, #3 tipus de document i unitat, #4 títol
\newcommand{\capcalera}[4]{%
  \gdef\fm@cursnom{#2}%
  \par\noindent
  \begin{tcolorbox}[colback=curs,colframe=curs,boxrule=0pt,arc=2.2mm,width=3.5cm,height=2.15cm,
      valign=center,left=3mm,right=2mm,top=1mm,bottom=1mm,nobeforeafter,box align=center]
    \color{curstext}{\footnotesize #1}\par\vspace{1pt}{\fontsize{19}{21}\selectfont\bfseries #2\par}%
  \end{tcolorbox}\hspace{5mm}%
  \begin{minipage}[c]{\dimexpr\linewidth-3.5cm-5mm-2.9cm\relax}
    {\color{cursfosc}\small\bfseries\MakeUppercase{#3}}\par\vspace{2pt}
    {\color{navy}\fontsize{20}{23}\selectfont\bfseries #4\par}
  \end{minipage}\hfill
  \begin{minipage}[c]{2.6cm}\raggedleft\includegraphics[width=2.6cm]{\fm@logo}\end{minipage}\par
  \vspace{2pt}}

% ---------- barres de secció del color del curs ----------
\renewcommand{\seccio}[2]{\par\needspace{\fmneed\baselineskip}\bigskip
  \begin{tcolorbox}[colback=curs,colframe=curs,coltext=curstext,arc=2mm,boxrule=0pt,left=3mm,right=3mm,top=1.2mm,bottom=1.2mm]
  {\large\bfseries #1\quad #2}
  \end{tcolorbox}\vspace{-2pt}}

% ---------- estrelles de dificultat del color del curs (a Batxillerat no se'n posen) ----------
\@ifundefined{estrelles}{}{%
  \renewcommand{\estrelles}[1]{{\color{curs}\ifcase#1\or\faStar\faStar[regular]\faStar[regular]\or\faStar\faStar\faStar[regular]\or\faStar\faStar\faStar\fi}}}

% ---------- capçalera i peu de les pàgines ----------
% primera pàgina: logos institucionals al peu
\fancypagestyle{primera}{\fancyhf{}\renewcommand{\headrulewidth}{0pt}
  \fancyfoot[L]{\raisebox{-0.32\height}{\includegraphics[height=0.78cm]{logo-gva.png}}\hspace{5mm}%
    \raisebox{-0.32\height}{\includegraphics[height=0.78cm]{logo-vermellar.png}}}
  \fancyfoot[R]{\scriptsize\color{muted}\href{https://creativecommons.org/licenses/by-nc-sa/4.0/deed.ca}{\ccbyncsa}\ \ CC BY-NC-SA 4.0 · fentmates.org\qquad\small\thepage}}
% la resta: el símbol a l'esquerra i el curs a la dreta, amb el títol al centre com abans
\fancyhead[L]{\raisebox{-0.3\height}{\includegraphics[height=0.62cm]{\fm@marca}}}
\fancyhead[R]{\small\bfseries\color{cursfosc}\fm@cursnom}
\renewcommand{\headrule}{\color{curs!45}\hrule height \headrulewidth}
\makeatother

%   \fmcurs{1eso}
% i, al començament del document:
%   \capcalera{Matemàtiques}{1r ESO}{Fitxa d'exercicis · Unitat 2}{Divisibilitat}
% Opcional abans de \input: \def\fmneed{14} (espai mínim abans de cada secció).
% =====================================================================
\makeatletter
\def\fm@logodir{../../logo/pdf}
\providecommand{\fmneed}{10}
\def\fm@cursnom{}

% clau del curs, color i color del text damunt del color
\newcommand{\fm@curs}[3]{\@namedef{fm@color@#1}{#2}\@namedef{fm@text@#1}{#3}}
\fm@curs{1eso}{E6B84A}{navy}
\fm@curs{2eso}{3B82F6}{white}
\fm@curs{3eso}{FF7A50}{white}
\fm@curs{4eso-a}{A78BFA}{white}
\fm@curs{4eso-b}{6D28D9}{white}
\fm@curs{1batx-mi}{9B1C1C}{white}
\fm@curs{1batx-ccss}{EF6B6B}{white}
\fm@curs{2batx-mii}{1E6B47}{white}
\fm@curs{2batx-ccss}{2BB3A3}{white}

\newcommand{\fmcurs}[1]{%
  \@ifundefined{fm@color@#1}{\PackageError{fentmates}{Curs desconegut: #1}{}}{}%
  \edef\fm@tmp{\noexpand\definecolor{curs}{HTML}{\@nameuse{fm@color@#1}}%
               \noexpand\colorlet{curstext}{\@nameuse{fm@text@#1}}}\fm@tmp
  \colorlet{cursfosc}{curs!72!black}%
  \def\fm@logo{\fm@logodir/vertical-#1.pdf}%
  \def\fm@marca{\fm@logodir/marca-#1.pdf}}

% ---------- capçalera de la primera pàgina ----------
% #1 matèria (petit), #2 curs, #3 tipus de document i unitat, #4 títol
\newcommand{\capcalera}[4]{%
  \gdef\fm@cursnom{#2}%
  \par\noindent
  \begin{tcolorbox}[colback=curs,colframe=curs,boxrule=0pt,arc=2.2mm,width=3.5cm,height=2.15cm,
      valign=center,left=3mm,right=2mm,top=1mm,bottom=1mm,nobeforeafter,box align=center]
    \color{curstext}{\footnotesize #1}\par\vspace{1pt}{\fontsize{19}{21}\selectfont\bfseries #2\par}%
  \end{tcolorbox}\hspace{5mm}%
  \begin{minipage}[c]{\dimexpr\linewidth-3.5cm-5mm-2.9cm\relax}
    {\color{cursfosc}\small\bfseries\MakeUppercase{#3}}\par\vspace{2pt}
    {\color{navy}\fontsize{20}{23}\selectfont\bfseries #4\par}
  \end{minipage}\hfill
  \begin{minipage}[c]{2.6cm}\raggedleft\includegraphics[width=2.6cm]{\fm@logo}\end{minipage}\par
  \vspace{2pt}}

% ---------- barres de secció del color del curs ----------
\renewcommand{\seccio}[2]{\par\needspace{\fmneed\baselineskip}\bigskip
  \begin{tcolorbox}[colback=curs,colframe=curs,coltext=curstext,arc=2mm,boxrule=0pt,left=3mm,right=3mm,top=1.2mm,bottom=1.2mm]
  {\large\bfseries #1\quad #2}
  \end{tcolorbox}\vspace{-2pt}}

% ---------- estrelles de dificultat del color del curs (a Batxillerat no se'n posen) ----------
\@ifundefined{estrelles}{}{%
  \renewcommand{\estrelles}[1]{{\color{curs}\ifcase#1\or\faStar\faStar[regular]\faStar[regular]\or\faStar\faStar\faStar[regular]\or\faStar\faStar\faStar\fi}}}

% ---------- capçalera i peu de les pàgines ----------
% primera pàgina: logos institucionals al peu
\fancypagestyle{primera}{\fancyhf{}\renewcommand{\headrulewidth}{0pt}
  \fancyfoot[L]{\raisebox{-0.32\height}{\includegraphics[height=0.78cm]{logo-gva.png}}\hspace{5mm}%
    \raisebox{-0.32\height}{\includegraphics[height=0.78cm]{logo-vermellar.png}}}
  \fancyfoot[R]{\scriptsize\color{muted}\href{https://creativecommons.org/licenses/by-nc-sa/4.0/deed.ca}{\ccbyncsa}\ \ CC BY-NC-SA 4.0 · fentmates.org\qquad\small\thepage}}
% la resta: el símbol a l'esquerra i el curs a la dreta, amb el títol al centre com abans
\fancyhead[L]{\raisebox{-0.3\height}{\includegraphics[height=0.62cm]{\fm@marca}}}
\fancyhead[R]{\small\bfseries\color{cursfosc}\fm@cursnom}
\renewcommand{\headrule}{\color{curs!45}\hrule height \headrulewidth}
\makeatother

\fmcurs{2eso}

\begin{document}
\thispagestyle{primera}

\capcalera{Matemáticas}{2\textsuperscript{o} ESO}{Soluciones de la ficha · Unidad 1}{Números naturales y enteros}
\vspace{8pt}

\begin{tcolorbox}[colback=soft,colframe=soft,arc=2mm,left=3mm,right=3mm,top=1.5mm,bottom=1.5mm]
\small Intenta resolver cada ejercicio antes de mirar su solución. Si no te sale el mismo resultado, repasa los pasos o consulta los apuntes interactivos de la unidad.
\end{tcolorbox}

% =====================================================
\seccio{1}{Los números naturales}

\ex{1} \begin{tasks}(2)
  \task $247=12\cdot\res{20}+\res{7}$
  \task $1000=37\cdot\res{27}+\res{1}$
  \task $505=5\cdot\res{101}+\res{0}$ (exacta)
  \task $89=90\cdot\res{0}+\res{89}$
\end{tasks}

\ex{1} \begin{tasks}(1)
  \task $7-12=-5$: \res{no} es natural (es negativo).
  \task $15\cdot4=\res{60}$: sí.
  \task $20:6$ \res{no} es exacta ($20=6\cdot3+2$).
  \task $6:3+1=\res{3}$: sí.
  \task $3-3=\res{0}$: sí (el 0 es natural).
\end{tasks}

\ex{2} $D=15\cdot23+9=\res{354}$. El resto \res{no} puede ser 17, porque tiene que ser menor que el divisor (15).

\ex{2} \begin{tasks}(2)
  \task $7\cdot 99=7\cdot 100-7=\res{693}$
  \task $12\cdot 15+12\cdot 5=12\cdot 20=\res{240}$
  \task $35\cdot 102=35\cdot 100+35\cdot 2=\res{3570}$
  \task $8\cdot (10+5)=80+40=\res{120}$
\end{tasks}


% =====================================================
\seccio{2}{Múltiplos y divisores}

\ex{1} \begin{tasks}(3)
  \task $7,\ 14,\ 21,\ 28,\ 35$
  \task $12,\ 24,\ 36,\ 48,\ 60$
  \task $15,\ 30,\ 45,\ 60,\ 75$
\end{tasks}

\ex{1} \begin{tasks}(1)
  \task $\text{Div}(24)=\{\res{1,\ 2,\ 3,\ 4,\ 6,\ 8,\ 12,\ 24}\}$
  \task $\text{Div}(36)=\{\res{1,\ 2,\ 3,\ 4,\ 6,\ 9,\ 12,\ 18,\ 36}\}$
  \task $\text{Div}(45)=\{\res{1,\ 3,\ 5,\ 9,\ 15,\ 45}\}$
  \task $\text{Div}(60)=\{\res{1,\ 2,\ 3,\ 4,\ 5,\ 6,\ 10,\ 12,\ 15,\ 20,\ 30,\ 60}\}$
  \task $\text{Div}(17)=\{\res{1,\ 17}\}$
\end{tasks}

\ex{2} \begin{tasks}(1)
  \task \res{V}: $48=6\cdot8$.
  \task \res{V}: $77=7\cdot11$.
  \task \res{F}: 5 es \emph{divisor} de 15.
  \task \res{V}: $n=1\cdot n$.
  \task \res{V}: $0=9\cdot0$.
\end{tasks}


% =====================================================
\seccio{3}{Criterios de divisibilidad. Números primos}

\ex{2} \begin{tasks}(2)
  \task $1320$: \res{2, 3, 5, 10, 11}
  \task $4455$: \res{3, 5, 9, 11}
  \task $7125$: \res{3, 5}
  \task $2002$: \res{2, 11}
\end{tasks}

\ex{2} \begin{tasks}(1)
  \task $4+\square+7=11+\square\in\{12,15,18\}\Rightarrow\square=\res{1,\ 4,\ 7}$
  \task Tiene que acabar en 0: $\square=\res{0}$ (230).
  \task $5+\square+2=9\Rightarrow\square=\res{2}$ (522)
\end{tasks}

\ex{1} \begin{tasks}(2)
  \task $31$: \res{primo}
  \task $51$: \res{compuesto: $51=3\cdot 17$}
  \task $57$: \res{compuesto: $57=3\cdot 19$}
  \task $61$: \res{primo}
  \task $87$: \res{compuesto: $87=3\cdot 29$}
  \task $91$: \res{compuesto: $91=7\cdot 13$}
  \task $97$: \res{primo}
  \task $1$: \res{ni primo ni compuesto}
\end{tasks}


% =====================================================
\seccio{4}{Descomposición en factores primos}

\ex{1} \begin{tasks}(3)
  \task $60=\res{2^{2}\cdot 3\cdot 5}$
  \task $84=\res{2^{2}\cdot 3\cdot 7}$
  \task $126=\res{2\cdot 3^{2}\cdot 7}$
  \task $360=\res{2^{3}\cdot 3^{2}\cdot 5}$
  \task $1000=\res{2^{3}\cdot 5^{3}}$
  \task $1176=\res{2^{3}\cdot 3\cdot 7^{2}}$
\end{tasks}

\ex{2} \begin{tasks}(3)
  \task $\res{600}$
  \task $\res{108}$
  \task $\res{1155}$
\end{tasks}

\ex{3} $72=2^3\cdot3^2$. Para que todos los exponentes sean pares falta un 2: $72\cdot\res{2}=144=2^4\cdot3^2=\res{12}^2$.


% =====================================================
\seccio{5}{Mínimo común múltiplo y máximo común divisor}

\ex{1} \begin{tasks}(4)
  \task $\res{36}$
  \task $\res{6}$
  \task $\res{40}$
  \task $\res{4}$
\end{tasks}

\ex{2} \begin{tasks}(1)
  \task $72=2^{3}\cdot 3^{2},\ 60=2^{2}\cdot 3\cdot 5$: \quad $\text{mcm}=\res{360}$,\ \ $\text{MCD}=\res{12}$
  \task $84=2^{2}\cdot 3\cdot 7,\ 126=2\cdot 3^{2}\cdot 7$: \quad $\text{mcm}=\res{252}$,\ \ $\text{MCD}=\res{42}$
  \task $45=3^{2}\cdot 5,\ 60=2^{2}\cdot 3\cdot 5,\ 75=3\cdot 5^{2}$: \quad $\text{mcm}=\res{900}$,\ \ $\text{MCD}=\res{15}$
  \task $140=2^{2}\cdot 5\cdot 7,\ 300=2^{2}\cdot 3\cdot 5^{2}$: \quad $\text{mcm}=\res{2100}$,\ \ $\text{MCD}=\res{20}$
  \task $16=2^{4},\ 27=3^{3}$: \quad $\text{mcm}=\res{432}$,\ \ $\text{MCD}=\res{1}$
\end{tasks}

\ex{2} $\text{MCD}\cdot\text{mcm}=12\cdot360=\res{4320}$ \quad $72\cdot60=\res{4320}$ \ \checkmark

\ex{3} $18\cdot b=6\cdot90=540\Rightarrow b=540:18=\res{30}$. Comprobación: $\text{MCD}(18,30)=6$ y $\text{mcm}(18,30)=90$.

\ex{2} \begin{tasks}(3)
  \task \res{Sí} (MCD $=1$)
  \task \res{No} (MCD $=7$)
  \task \res{Sí} (MCD $=1$)
\end{tasks}


% =====================================================
\seccio{6}{Problemas de mcm y MCD}

\ex{2} $\text{mcm}(15,20)=60$ min $\Rightarrow$ \res{8:00}

\ex{2} $\text{MCD}(60,84)=\res{12}$ ramos; $60:12=\res{5}$ rosas y $84:12=\res{7}$ claveles cada ramo.

\ex{2} $\text{MCD}(480,360)=\res{120\text{ cm}}$; \quad $4\cdot3=\res{12}$ baldosas.

\ex{2} $\text{mcm}(8,12,18)=2^3\cdot3^2=\res{72\text{ s}}$

\ex{2} $\text{MCD}(150,210)=\res{30}$ libros; $5+7=\res{12}$ paquetes.

\ex{3} $\text{mcm}(6,8,10)=120$; \quad $120+3=\res{123}$ caramelos.

\ex{3} $\text{mcm}(24,36)=72$ dientes: $72:24=\res{3}$ vueltas la primera y $72:36=\res{2}$ la segunda.


% =====================================================
\seccio{7}{Los números enteros}

\ex{1} \begin{tasks}(6)
  \task $\res{-3}$
  \task $\res{-2}$
  \task $\res{-50}$
  \task $\res{+1200}$
  \task $\res{-44}$
  \task $\res{+15}$
\end{tasks}

\ex{1} \begin{tasks}(2)
  \task $\lvert -7\rvert =\res{7}$, \ $\text{op}(-7)=\res{7}$
  \task $\lvert 12\rvert =\res{12}$, \ $\text{op}(12)=\res{-12}$
  \task $\lvert 0\rvert =\res{0}$, \ $\text{op}(0)=\res{0}$
  \task $\lvert -25\rvert =\res{25}$, \ $\text{op}(-25)=\res{25}$
\end{tasks}

\ex{1} $\res{-12<-8<-3<-1<0<2<5<7}$

\ex{2} \begin{tasks}(2)
  \task $\res{-3,{-2},{-1},0,1,2}$
  \task $\res{-2,{-1},0,1,2}$
  \task $\res{-5,\ 5}$
  \task $\res{-2}$
\end{tasks}

\ex{2} El \res{martes}. $-5<-3<-2<-1<0<1<4$. \quad $4-(-5)=4+5=\res{9}$\,°C


% =====================================================
\seccio{8}{Operaciones con enteros}

\ex{1} \begin{tasks}(4)
  \task $(-7)+(-5)=\res{-12}$
  \task $(-9)+4=\res{-5}$
  \task $12+(-15)=\res{-3}$
  \task $(-6)-(-10)=\res{4}$
  \task $8-(+13)=\res{-5}$
  \task $(-3)-5=\res{-8}$
  \task $0-(-7)=\res{7}$
  \task $(-11)+11=\res{0}$
\end{tasks}

\ex{1} \begin{tasks}(2)
  \task $5-8+3-7+2=(5+3+2)-(8+7)=10-15=\res{-5}$
  \task $-4-6+10-3=10-(4+6+3)=10-13=\res{-3}$
  \task $12-15-9+20=(12+20)-(15+9)=32-24=\res{8}$
  \task $-1+2-3+4-5+6=(2+4+6)-(1+3+5)=12-9=\res{3}$
\end{tasks}

\ex{2} \begin{tasks}(1)
  \task $7-(3-5+2)=7-3+5-2=\res{7}$
  \task $-(4-9)+(-2+6)=5+4=\res{9}$
  \task $10-\left[3-(5-8)\right]=10-\left[3-(-3)\right]=10-6=\res{4}$
  \task $-\left[-(2-7)-(-3+1)\right]=-\left[5-(-2)\right]=-7=\res{-7}$
\end{tasks}

\ex{1} \begin{tasks}(2)
  \task $(-6)\cdot (-4)=\res{24}$
  \task $7\cdot (-8)=\res{-56}$
  \task $(-45):9=\res{-5}$
  \task $(-56):(-7)=\res{8}$
  \task $0:(-5)=\res{0}$
  \task $(-2)\cdot (-3)\cdot (-5)=\res{-30}$
  \task $(-1)\cdot (-1)\cdot (-1)\cdot (-1)=\res{1}$
  \task $72:(-8)=\res{-9}$
\end{tasks}

\ex{2} \begin{tasks}(3)
  \task $\square=\res{-7}$
  \task $\square=\res{-30}$
  \task $\square=\res{9}$
  \task $\square=\res{4}$
  \task $\square=\res{-7}$
  \task $\square=\res{5}$
\end{tasks}

\ex{2} $-120+45-80=\res{-155}$: a \res{155 m} de profundidad.

\ex{2} $-495-(-570)=-495+570=\res{75}$ años.

\ex{3} $230-3\cdot 85+120=230-255+120=\res{95}$\,€. Sí: después de los recibos, $230-255=\res{-25}$\,€.


% =====================================================
\seccio{9}{Operaciones combinadas}

\ex{1} \begin{tasks}(2)
  \task $5-3\cdot 4=5-12=\res{-7}$
  \task $(5-3)\cdot 4=2\cdot 4=\res{8}$
  \task $18:(-3)+2\cdot (-5)=-6-10=\res{-16}$
  \task $-20:4\cdot (-2)=-5\cdot (-2)=\res{10}$
  \task $3\cdot (-2)-4\cdot (-3)=-6+12=\res{6}$
\end{tasks}

\ex{2} \begin{tasks}(1)
  \task $4-2\cdot \left[3-(6-10)\right]=4-2\cdot \left[3+4\right]=4-14=\res{-10}$
  \task $(-3)\cdot (2-7)+15:(-5)=(-3)\cdot (-5)-3=15-3=\res{12}$
  \task $\left[(-8)+5\right]\cdot \left[6-(-2)\right]=(-3)\cdot 8=\res{-24}$
  \task $-6-\left[4-3\cdot (2-5)\right]=-6-\left[4+9\right]=-6-13=\res{-19}$
  \task $25:(-5)-(-3)\cdot (4-6):2=-5-(-3)\cdot (-2):2=-5-6:2=-5-3=\res{-8}$
\end{tasks}

\ex{3} Ha restado antes de dividir y multiplicar. $12-4:2\cdot (-3)=12-2\cdot (-3)=12+6=\res{18}$

\ex{3} \begin{tasks}(3)
  \task $(6-4)\cdot (3+2)=\res{10}$
  \task $6-4\cdot (3+2)=\res{-14}$
  \task $-8:(2-6)=\res{2}$
\end{tasks}


% =====================================================
\seccio{10}{Potencias y raíces con enteros}

\ex{1} \begin{tasks}(4)
  \task $(-3)^{4}=\res{81}$
  \task $-3^{4}=\res{-81}$
  \task $(-2)^{5}=\res{-32}$
  \task $(-1)^{100}=\res{1}$
  \task $(-1)^{25}=\res{-1}$
  \task $(-10)^{3}=\res{-1000}$
  \task $0^{7}=\res{0}$
  \task $(-7)^{0}=\res{1}$
\end{tasks}

\ex{2} \begin{tasks}(3)
  \task $(-2)^3\cdot(-2)^5=\res{(-2)^{8}}$
  \task $(-5)^9:(-5)^6=\res{(-5)^{3}}$
  \task $\left[(-3)^2\right]^4=\res{(-3)^{8}}$
  \task $2^5\cdot(-3)^5=\res{(-6)^{5}}$
  \task $(-12)^4:3^4=\res{(-4)^{4}}$
  \task $7^3\cdot7^4:7^5=\res{7^{2}}$
\end{tasks}

\ex{2} \begin{tasks}(2)
  \task $(-2)^7:(-2)^4=(-2)^3=\res{-8}$
  \task $\left[(-1)^3\right]^5\cdot(-1)^2=(-1)^{17}=\res{-1}$
  \task $5^4\cdot2^4=10^4=\res{10000}$
  \task $(-6)^3:(-3)^3=2^3=\res{8}$
\end{tasks}

\ex{1} \begin{tasks}(3)
  \task $\res{\pm9}$
  \task $\res{\pm12}$
  \task \res{No existe}
  \task $\res{0}$
  \task $\res{\pm20}$
  \task \res{No existe}
\end{tasks}

\ex{2} \begin{tasks}(2)
  \task $\sqrt{50}\approx\res{7}$, resto $\res{1}$: \ $50=7^2+1$
  \task $\sqrt{90}\approx\res{9}$, resto $\res{9}$: \ $90=9^2+9$
  \task $\sqrt{130}\approx\res{11}$, resto $\res{9}$: \ $130=11^2+9$
  \task $\sqrt{200}\approx\res{14}$, resto $\res{4}$: \ $200=14^2+4$
\end{tasks}

\ex{2} \begin{tasks}(2)
  \task $(-2)^{3}+3^{2}=-8+9=\res{1}$
  \task $-2^{4}+(-2)^{4}=-16+16=\res{0}$
  \task $5-2\cdot \sqrt{36}+(-1)^{7}=5-2\cdot 6-1=5-12-1=\res{-8}$
  \task $(-3)^{2}\cdot 2-\sqrt{49}:(-7)=9\cdot 2-7:(-7)=18-(-1)=\res{19}$
  \task $\left[(-2)^{2}-3^{2}\right]^{2}:(-5)=\left[4-9\right]^{2}:(-5)=25:(-5)=\res{-5}$
  \task $\sqrt{25-9}-(-3)^{3}:9=\sqrt{16}-(-27):9=4-(-3)=\res{7}$
\end{tasks}

\ex{2} \begin{tasks}(5)
  \task $\res{+}$
  \task $\res{-}$
  \task $\res{-}$
  \task $\res{-}$
  \task $\res{+}$
\end{tasks}

\ex{3} $1225=5^2\cdot7^2\Rightarrow\sqrt{1225}=5\cdot7=\res{35\text{ m}}$; \quad $4\cdot35=\res{140\text{ m}}$ de valla.


\end{document}
